\section{Multi-Objective Policy Synthesis}

We formalize the policy synthesis problem for a fixed number of goals, and then introduce our auction-based framework for this setting. 
Afterwards, we will describe how our framework naturally extends to the setting with dynamic goal appearance and disappearance.

\textbf{Centralized optimal policy synthesis for multiple objectives.}
We use \textit{multi-objective Markov decision processes} (\momdps) to model uncertain environments with multiple reward objectives.
An \momdp with \(m\in \N_{>0}\) objectives is a tuple \(\mathcal{M}=(S,A,T,\mu,R,H)\), where \(S\) is the \textit{state} space, \(A\) is the \textit{action} space, \(T\colon S\times A\to \dist(S)\) is the \textit{transition} kernel, \(\mu\in \dist(S)\) is the \textit{initial state} distribution, \(R=\{r_i\colon S\times A\to \R\}_{i\in [1;m]}\) is the set of \textit{reward functions}, and $H\in \mathbb{N}_{>0}$ is the \textit{time horizon}.
A (Markov) \textit{policy} is a sequence \(\pi=\{\pi_t\}_{t=0}^{H-1}\), where each \(\pi_t:S\to \dist(A)\) maps the current state to a distribution over actions.
For a trace \(\rho=s^0a^0s^1a^1\ldots s^H\), the accumulated \textit{local} reward of objective \(i\) is \(w_i(\rho)=\sum_{t=0}^{H-1} r_i(s^t,a^t)\), and the \textit{global} reward is \(w(\rho)=\sum_{i=1}^m w_i(\rho)\).
The \textit{value} of a policy is \(V^\pi=\E^\pi[w(\rho)]\), and an \textit{optimal} policy is any policy \(\pi^*\in \arg\max_\pi V^\pi\).

\smallskip\noindent
\textbf{Lower complexity bound.}\todo{this part needs work}
We establish that the optimal policy synthesis problem is NP-hard, and we show this already on a simplified setting:
	Suppose the transition function $T$ of the \momdp is deterministic, the goals are stationary such that the location of the $i$-th goal is always $g_i\in S$, and for any policy, the minimum travel time from a given location $x\in S$ to the $i$-th goal is $t_i = d(x,g_i)$. 
	(Alternatively, $t_i$ can be interpreted as the length of the shortest path from $x$ to $g_i$ when the agent moves at the unit velocity.)

\begin{theorem}\label{thm:np hardness}
	Finding the optimal policy for the \momdp is NP-hard.
\end{theorem}

\begin{proof}
	The simpler setting with deterministic agents and stationary goals subsumes the Euclidean traveling salesman problem, where the deadlines are all infinity, and the goal is to find the shortest path passing through all the goals.
	This problem is known to be NP-complete~\cite{DBLP:journals/tcs/Papadimitriou77}, and our deadline-restricted case is at least as hard as it.
	\KM{There is a slight difference, that we allow revisiting a state. However, I found on discussion forums that this does not change the complexity, find a reference. Also see metric traveling salesman, which is more general (we could then replace the Euclidean distance requirement in the \momdp state space with the requirement of just having a metric.}
\end{proof}

The NP-hardness of the optimal policy synthesis problem motivates us to find heuristics to solve the problem.
Moreover, our heuristic must seamlessly support dynamic addition or deletion of objectives (modeled as reward functions), which is our ultimate goal.
Towards our goal, we identify a key structural property of the optimal policy below.

\begin{proposition}\label{prop:scheduling is the optimal choice}
	Consider the simpler case of deterministic agent with stationary goals.
	There exists a permutation $\sigma^*\in \Sigma_m$, where $\Sigma_m$ is the set of all permutations of $[1;m]$, such that the optimal policy $\pi^*$ pursues the goals in the order $\sigma^*(1)\ldots\sigma^*(m)$, and $\pi^*$ takes the shortest paths from the initial agent location to goal $\sigma^*(1)$ and between $\sigma^*(i)$ and $\sigma^*(i+1)$, for every $i\in [1;m-1]$.
\end{proposition}
The above proposition effectively reduces the problem to an optimal \emph{scheduling} problem.
Furthermore, Theorem~\ref{thm:np hardness} suggests that a polynomial-time centralized optimal algorithm does not exist, unless P=NP, already for the simpler deterministic case with stationary goals.
This motivates us to seek for a heuristic approach that is aimed towards finding the optimal schedule over the goals whose existence is guaranteed by Proposition~\ref{prop:scheduling is the optimal choice}.

A natural deterministic baseline is the \emph{least slack time first} (LSTF) rule, which schedules the goal with smallest slack $\slack_i(x,t)\coloneqq d_i-t-d(x,g_i)$.
However, this notion of urgency is tailored to stationary goals and deterministic travel times, and does not align with our full \momdp objective once targets move stochastically or may appear and disappear over time.
For this reason, we use a more general benchmark that we call \emph{reward-aware urgency first} (RAUF).
\KM{Is this used in the literature? If yes, some citation would be nice.}
At every decision point, RAUF selects one of the goals with maximal \emph{reward-aware urgency}, namely the marginal continuation-value gain from granting that goal immediate control rather than deferring control to competing goals; ties are broken uniformly at random.
Unlike slack, this urgency notion directly accounts for discounted rewards, missed-deadline penalties, stochastic target motion, and future target arrivals or removals.
Our auction mechanism is designed to approximate this reward-aware urgency ordering in a decentralized manner while remaining tractable in settings where computing the globally optimal schedule is infeasible.
\GS{Also talk about how urgency encodes agent cooperation.}


\section{Auction-Based Multi-Policy Framework}
We consider the compositional approach to policy synthesis for \momdps, where we will design a selfish, \emph{local} policy maximizing each individual value component, towards the fulfillment of some required global coordination requirements.
The main crux is in the composition process, where each local policy may propose a different action, but the composition must decide one of the actions that will be actually executed.
Importantly, the composition must be implementable in a distributed manner, meaning we will \emph{not} use any global policy that would pick an action by analyzing all local policies and their reward functions.

We present a novel \textit{auction}-based RL framework for compositional policy synthesis for \momdps.
In our framework, not only do the local policies emit actions, but they also submit \emph{bids} within a given bound $\beta\in \mathbb{N}_{>0}$ to win the privilege of executing their actions for a given number of time steps $\tau\in \mathbb{N}_{>0}$.
We consider two models of bidding, called \poorman and \allpay.
In the \poorman setting, the bids are non-negative \textit{integers} in the interval $[0;\beta]$, and in the \allpay setting, the bids are non-negative \textit{real} numbers in the interval $[0,\beta]$.%
\footnote{The distinction in bid types---\textit{integers} for \poorman and \textit{reals} for \allpay---is purely technical. 
The synthesis of local policies will involve a game-theoretic analysis involving Nash equilibria.
For real-valued bids in \poorman, it is unclear if Nash equilibria would always exist, owing to the continuous bid values and discontinuities in the payoff profile due to the tie-breaking mechanism~\cite{simon1990discontinuous}. 
For integer-valued bids in \allpay, the shape of the Nash equilibria policies is only known for the two-policy case~\cite{dziubinski2023discrete}, but the general case is still open.
We leave these unresolved cases for a future work.
}
In both settings, the highest bidder's actions get executed for the following $\tau$ consecutive steps, with bidding ties being resolved uniformly at random.
To encourage truthfulness, we penalize the policies proportional to their bids scaled with a constant $\rho\in (0,1)$, and which policies receive penalties differ between the two settings:
for \poorman, only the executed policy (one of the highest bidders) pays the penalty, whereas for \allpay, all policies pay penalties.
Through this mechanism, policies select their bids to reflect how critical their action is in the current state, and the bidding penalty incentivizes truthful urgency.
This way, we obtain a purely decentralized scheme to coordinate local policies in a given \momdp.


The parameters $\tau$, $\beta$, and $\rho$ denote the \emph{bidding interval}, \emph{maximum bid}, and \emph{bid penalty coefficient}, respectively. Ablation studies in Section~\ref{sec:experimental-results} highlight their importance. If $\tau$ is too small, policies switch too frequently, risking suboptimal outcomes. If $\beta$ is too small, policies cannot express preference strength, leading to near-random control. If $\rho$ is too large, policies tend to abstain from bidding.

To connect the bidding urgencies with the objectives, we formulate the local policy synthesis question as a game theoretic problem defined below.

\begin{definition}[Markov bidding game]\label{def:markov-bidding-game}
	Let \(\mathcal{M}=(S,A,T,\mu,R,H)\) be an \momdp with \(m\) objectives, let \(\tau\in\N_{>0}\) be the bidding interval, let \(\beta\in\R_{>0}\) be the maximum bid, let \(\rho\in(0,1)\) be the bid penalty coefficient, and let \(\mathsf{P}\in\{\poorman,\allpay\}\) be the penalty model.
	The induced Markov bidding game between the $m$ local policies is \(\mathcal{G}^{\mathcal{M}}_{\tau,\beta,\rho,\mathsf{P}}=(\hat S,\hat A,\hat T,\hat R,\hat\mu,H)\), where:
	\begin{itemize}
		\item \(\hat S=S\times [1;m]\times [0;\tau-1]\). A state \((s,k,q)\) records the \momdp state \(s\), the current controlling policy \(k\), and the number \(q\) of steps remaining before the next bidding round. When \(q=0\), a new bidding round occurs and the current controller \(k\) is ignored.
		\item \(\hat A=(A\times [0;\beta])^m\) for $\mathsf{P}=\poorman$ and $\hat A=(A\times [0,\beta]^m)$ for $\mathsf{P}=\allpay$. Each policy \(i\) independently chooses an action \(a_i\in A\) and a bid $b_i$, where \(b_i\in[0;\beta]\) or $b_i\in [0,\beta]$ depending on $\mathsf{P}$. We write a joint action as \((\vec a,\vec b)=(a_i,b_i)_{i\in[1;m]}\).
		\item Let \(W(\vec b)=\{i\in[1;m]\mid b_i=\max_j b_j\}\) be the set of highest bidders. The transition kernel \(\hat T:\hat S\times\hat A\to\dist(\hat S)\) is defined by
		\begin{equation}
			\hat T((s,k,q),(\vec a,\vec b),(s',k',q'))=
			\begin{cases}
				T(s,a_k,s') & q>0 \land k'=k \land q'=q-1,\\
				\frac{1}{\abs{W(\vec b)}}T(s,a_{k'},s') & q=0 \land k'\in W(\vec b)\land q'=\tau-1,\\
				0 & \text{otherwise}.
			\end{cases}
		\end{equation}
		Intuitively, the case ``$q>0$'' means it is not yet a bidding round, and the currently active policy $k$ continues to be used, and the case ``$q=0$'' means it is a bidding round, and a new policy is selected via tie-breaking over the set $W$.
		\item For each policy \(i\), the reward is inherited from the \momdp reward \(r_i\), with a bid penalty charged at bidding rounds:
		\begin{equation}
			\hat r_i((s,k,q),(\vec a,\vec b),(s',k',q'))
			=
			r_i(s,a_{k'})
			-
			\rho b_i\cdot \mathbf{1}\!\left[q=0\land (i=k'\lor \mathsf{P}=\allpay)\right].
		\end{equation}
		\item The initial distribution \(\hat\mu\) satisfies \(\hat\mu(s,k,0)=\mu(s)/m\) for every \(k\in[1;m]\), and \(\hat\mu(s,k,q)=0\) for \(q>0\).
	\end{itemize}
\end{definition}
Thus a local policy $i$ in the \momdp $\M$ is a strategy $\pi_i\colon \hat S\to \hat A_i$ in the Markov bidding game $\mathcal{G}^{\mathcal{M}}_{\tau,\beta,\rho,\mathsf{P}}$ maximizing its own expected cumulative values of $\mathbb{E}[\sum_{t=0}^{H-1}\hat r_i^t]$, against every other strategy $\pi_j$ that is maximizing its own expected cumulative reward $\mathbb{E}[\sum_{t=0}^{H-1}\hat r_j^t]$, while the auction rule determines which local action is executed.
We write $\pi_i$ as a pair $(\pi_i^a,\pi_i^b)$, where \(\pi_i^a\colon \hat S\to A\) is the action component and \(\pi_i^b\colon \hat S\to [0;\beta]\) is the bid component  for \poorman and $\pi_i^b\colon \hat S\to [0,\beta]$ is the same for \allpay.
We next analyze the bidding incentives at a single decision state, using continuation values from this game.


%\smallskip
%\noindent\textbf{Existence of Nash equilibria.}
%Continuous bidding makes the Markov bidding game have discontinuous payoffs at bid ties, since an arbitrarily small increase in bid by any of the tied policies will change the winning policy.
%Simon and Zame~\citep{simon1990discontinuous} prove the existence of Nash equilibria for games with discontinuous payoffs by allowing the tie-breaking mechanism (formalized as ``endogenous payoff sharing rule''  at discontinuity points) to be computed as part of the equilibrium object.
%This does not directly imply that a Nash equilibrium would exist for any arbitrary tie-breaking mechanism fixed apriori, which is the case for us. 
%Nevertheless, we chose uniform tie-breaking due to its simplicity, and since the event that a tie happens has Lebesgue measure zero, in practice, 
%
%in the continuous all-pay equilibrium characterized below, bid ties occur with probability zero, so the particular tie-breaking rule is payoff-irrelevant for that equilibrium.
%The winner-pays result below should instead be read conditionally on the existence or selection of a pure equilibrium.

\smallskip
\noindent\textbf{Equilibrium continuation values.}
We consider local policies \(\pi_i=(\pi_i^a,\pi_i^b)\) that form a Nash equilibrium of the Markov bidding game $\mathcal{G}^{\mathcal{M}}_{\tau,\beta,\rho,\mathsf{P}}$.
For simplicity, assume $H$ is a multiple of $\tau$ (otherwise, extend the horizon to the next multiple of $\tau$).
The policies bid at \textit{bidding rounds} $t\in \{0,\tau,2\tau,\ldots,H-\tau\}$.
Given a policy index $i\in [1;m]$, time $t\leq H$, and state $s\in \hat S$ of $\mathcal{G}^{\mathcal{M}}_{\tau,\beta,\rho,\mathsf{P}}$, we write $V^*_{i,t}$ to denote the value policy $i$ achieves in the rest of the game starting at time $t$ and at state $s$, when all the policies follow a Nash equilibrium.
Since the horizon $H$ is finite, we can define $V^*_{i,t}$ using backward induction over $t$ starting from $t=H$.
Set \(V^*_{i,H}(s)=0\) for every \(i\) and \(s\).
Suppose the game is at state \(s\) and bidding round $t\in \{0,\tau,2\tau,\ldots,H-\tau\}$.

The \textit{winning continuation value} of a policy $i$ is the optimal value it can achieve upon winning bidding:
\begin{equation}
	\widetilde V^{\win}_{i,t}(s)
	\coloneqq
	\max_{\pi_i^a}
	\E^{\pi_i^a}\left[
		\sum_{q=0}^{\tau-1} r_i(s^{t+q},a^{t+q})
		+V^*_{i,t+\tau}(s^{t+\tau})
	\right],
\end{equation}
where the action sequence is generated by policy \(i\) during the time window $[t;t+\tau]$.
Similarly, the \textit{losing continuation value} is the least value it achieves upon losing the bidding:
\begin{equation}\label{eq:losing continuation value}
	\widetilde V^{\lose}_{i,t}(s)
	\coloneqq
	\min_{j\neq i}
	\E^{\pi_j^a}\left[
		\sum_{q=0}^{\tau-1} r_i(s^{t+q},a^{t+q})
		+V^*_{i,t+\tau}(s^{t+\tau})
	\right],
\end{equation}
where the action sequence during $[t;t+\tau]$ is generated by a different policy \(j\) using an action policy $\pi_j^a$ that maximizes $j$'s own winning continuation value.
The \textit{control gain} of policy \(i\) at \((s,t)\) is
\begin{equation}
	G_{i,t}(s)
	\coloneqq
	\widetilde V^{\win}_{i,t}(s)-\widetilde V^{\lose}_{i,t}(s).
\end{equation}
This quantity measures the advantage of winning control for the next bidding interval rather than losing it.
Given bids \(b=(b_1,\ldots,b_m)\), let
$W(b)\coloneqq \{k\in [1;m]\mid b_k=\max_{\ell\in[1;m]}b_\ell\}$ 
be the set of highest bidders, and let \(p_i(b)=\mathbf{1}[i\in W(b)]/\abs{W(b)}\) be the probability that \(i\) wins after uniform tie-breaking.
For a penalty model \(\mathsf{P}\in\{\poorman,\allpay\}\), the \textit{one-round bidding payoff} is
\begin{equation}
	U^{\mathsf{P}}_{i,t}(s,b)
	=
	\begin{cases}
		\widetilde V^{\lose}_{i,t}(s)+p_i(b)\left(G_{i,t}(s)-\rho b_i\right) & \mathsf{P}=\poorman,\\
		\widetilde V^{\lose}_{i,t}(s)+p_i(b)G_{i,t}(s)-\rho b_i & \mathsf{P}=\allpay.
	\end{cases}
\end{equation}
In the \poorman case, the bid cost is paid only in expectation over winning; in the \allpay case, the bid cost is paid regardless of the winner.
Effectively, in the Markov game, the goal of policy $i$ boils down to maximizing the one-round bidding payoff $U^{\mathsf{P}}_{i,t}(s,b)$ at each state $s$ and bidding round $t$. 

A mixed bidding policy for player \(i\) at \((s,t)\) is the bid component \(\pi^b_{i,t}(\cdot\mid s)\in\dist([0,\beta])\), and a mixed bidding profile is \(\pi^b_t(\cdot\mid s)=(\pi^b_{1,t}(\cdot\mid s),\ldots,\pi^b_{m,t}(\cdot\mid s))\).
We write \(b\sim\pi^b_t(\cdot\mid s)\) when bids are sampled independently as \(b_i\sim\pi^b_{i,t}(\cdot\mid s)\).
A mixed bidding profile \(\pi^{b*}_t(\cdot\mid s)\) is a Nash equilibrium of the one-round bidding game at \((s,t)\) if, for every player \(i\) and every alternative bid distribution $\pi^b_{i,t}$, where \(\pi^b_{i,t}(\cdot\mid s)\in\dist([0;\beta])\) for $\mathsf{P}=\poorman$ and $\mathsf{P}=\allpay$,
\begin{equation}
	\E_{b\sim\pi^{b*}_t(\cdot\mid s)}\!\left[U^{\mathsf{P}}_{i,t}(s,b)\right]
	\geq
	\E_{b_i\sim\pi^b_{i,t}(\cdot\mid s),\; b_{-i}\sim\pi^{b*}_{-i,t}(\cdot\mid s)}\!\left[U^{\mathsf{P}}_{i,t}(s,b_i,b_{-i})\right].
\end{equation}
For the \poorman setting, the existence of the Nash equilibrium is guaranteed due to the finiteness of the action space $\hat A$~\citep{nash1951noncooperative}, and for the \allpay setting, the same follows from known results in the literature~\citep{baye1996all}.
The equilibrium continuation value is then
\begin{equation}
	V^*_{i,t}(s)
	=
	\E_{b\sim \pi^{b*}_t(\cdot\mid s)}\left[U^{\mathsf{P}}_{i,t}(s,b)\right],
\end{equation}
which completes the backward-induction definition.

\smallskip
\noindent\textbf{Theoretical analysis of the bidding schemes.}
We will now prove that, if all policies follow the Nash equilibrium and if some mild assumptions hold true, then the policy with the highest control gain is guaranteed to get the control at each state and time.
Consequently, we obtain a greedy solution to maximize the global reward (the sum of local cumulative rewards).
Further, from Eq.~\eqref{eq:losing continuation value} it follows that the control gains are actually environment-aware, i.e., if the goal of policy $i$ is aligned with every other policy (like for the cat feeder, all cats are crowded in one corner), then the control gain of $i$ would be automatically low.
We state our two results here but defer the proofs to~\Cref{app:poorman-control-gain-proof,app:allpay-control-gain-proof}.


\begin{theorem}[\poorman favors high control gain]\label{thm:poorman-favors-high-control-gain}
	In the \poorman setting, consider a single bidding round $t$ at the agent state $s\in S$.
  	Suppose the parameters $\rho$ and $\beta$ are so chosen that $\rho\beta > \max_{j\in [1;m]}G_{j,t}(s)$.
	Let $i^*$ be the policy with the largest control gain, such that the following margin condition is fulfilled:
	\begin{align}\label{eq:all pay:gross gain gap}
		G_{i^*,t}(s) - \max_{j\neq i^*}G_{j,t}(s) > \rho.
	\end{align}
	Then in any pure Nash equilibrium, policy $i^*$ wins control.
\end{theorem}

\begin{theorem}[\allpay favors high control gain]\label{thm:allpay-favors-high-control-gain}
	In the \allpay setting with $m\geq 3$, consider a single bidding round $t$ at the agent state $s\in S$.
	Let $i^*$ and $j^*$ be policies for which the following gross gain order holds:
	\begin{align}\label{eqn:allpay urgency gap}
		G_{i^*,t} > G_{j^*,t} > \max_{k\notin \{i^*,j^*\}}G_{k,t}.
	\end{align}
	Then there is no Nash equilibrium with pure strategies, although a unique Nash equilibrium exists with mixed strategies.
	In this Nash equilibrium, policy $i^*$ wins control with probability $1-\sfrac{G_{j^*,t}}{2G_{i^*,t}}$.	
\end{theorem}

While these results do not automatically imply equivalence to the centralized optimal solution, we obtain an environment-aware, greedy approximation.
This gives a scheduling perspective on the auction: each bidding round picks the objective to receive control based on its urgency.
In Appendix~\ref{app:lstf-analysis}, we prove that in deterministic environments with stationary targets and varying deadlines, a Nash equilibrium in our auction mechanism always selects critically urgent goals, if any.
This matches the least slack time first (LSTF) rule, a widely used scheduling heuristic~\citep{brown2020optimal}.


%
%\smallskip
%\noindent\textbf{Theoretical analysis of the bidding schemes.}
%The first result says that the winner-pays mechanism can implement the largest-gain policy in arbitrarily good pure approximate equilibria.
%Recall that a pure bid profile is a pure \(\varepsilon\)-Nash equilibrium if no policy can improve its payoff by more than \(\varepsilon\) through a unilateral bid deviation.
%
%\begin{theorem}[Winner-pays approximate implementation]
%	In the \poorman setting, consider a single bidding round at state \(s\) and time \(t\).
%	Let \(i^*\) be a policy with largest control gain, let \(j^*\) be a policy with second-largest control gain, and suppose
%	\begin{equation}
%		G_{i^*,t}(s)>G_{j^*,t}(s).
%	\end{equation}
%	For any \(\varepsilon>0\) such that \((G_{j^*,t}(s)+\varepsilon)/\rho\leq \beta\) and \(G_{i^*,t}(s)>G_{j^*,t}(s)+\varepsilon\), there is a pure \(\varepsilon\)-Nash equilibrium in which policy \(i^*\) wins control with probability one.
%\end{theorem}
%
%\begin{proof}
%	Consider the bid profile \(b_{j^*}=G_{j^*,t}(s)/\rho\), \(b_{i^*}=(G_{j^*,t}(s)+\varepsilon)/\rho\), and \(b_k\leq b_{j^*}\) for every \(k\notin\{i^*,j^*\}\).
%	Policy \(i^*\) wins uniquely.
%	No losing policy can gain by deviating to another losing bid.
%	If \(j^*\) deviates to a winning bid, it must bid at least \(b_{i^*}\), giving payoff at most
%	\begin{equation}
%		\widetilde V^{\lose}_{j^*,t}(s)+G_{j^*,t}(s)-\rho b_{i^*}
%		=
%		\widetilde V^{\lose}_{j^*,t}(s)-\varepsilon,
%	\end{equation}
%	which is worse than losing.
%	Every other losing policy has control gain at most \(G_{j^*,t}(s)\), so the same argument applies.
%	
%	Policy \(i^*\) cannot gain by raising its bid.
%	Its best deviation is to lower its bid while still winning, which can improve its payoff by at most \(\varepsilon\), since it must bid above \(b_{j^*}\) to remain the unique winner.
%	Deviating to a losing bid gives payoff \(\widetilde V^{\lose}_{i^*,t}(s)\), while the proposed bid gives
%	\begin{equation}
%		\widetilde V^{\lose}_{i^*,t}(s)+G_{i^*,t}(s)-G_{j^*,t}(s)-\varepsilon,
%	\end{equation}
%	which is strictly larger by assumption.
%	Thus no policy can improve by more than \(\varepsilon\), so the profile is a pure \(\varepsilon\)-Nash equilibrium.
%\end{proof}
%
%The all-pay mechanism admits a standard mixed-equilibrium construction in the continuous complete-information setting.
%
%\begin{theorem}[All-pay concentration]
%	In the \allpay setting, consider a single bidding round at state \(s\) and time \(t\), with continuous bid space \([0,\beta]\).
%	Let \(i^*\) and \(j^*\) be policies satisfying
%	\begin{equation}
%		G_{i^*,t}(s)
%		>
%		G_{j^*,t}(s)
%		>
%		\max_{k\notin\{i^*,j^*\}}G_{k,t}(s),
%	\end{equation}
%	and suppose \(\rho\beta\geq G_{j^*,t}(s)\).
%	Then the continuous all-pay bidding game has no pure Nash equilibrium, and there is a mixed Nash equilibrium given by the following bid distributions:
%	\begin{itemize}
%		\item policy \(i^*\) samples uniformly from \([0,G_{j^*,t}(s)/\rho]\);
%		\item policy \(j^*\) places an atom of mass
%		\begin{equation}
%			\frac{G_{i^*,t}(s)-G_{j^*,t}(s)}
%			{G_{i^*,t}(s)}
%		\end{equation}
%		at \(0\), and otherwise has density \(\rho/G_{i^*,t}(s)\) on \((0,G_{j^*,t}(s)/\rho]\);
%		\item every other policy bids \(0\) with probability one.
%	\end{itemize}
%	In this equilibrium, policy \(i^*\) wins control with probability
%	\begin{equation}
%		1-\frac{G_{j^*,t}(s)}
%		{2G_{i^*,t}(s)}.
%	\end{equation}
%\end{theorem}
%
%\begin{proof}
%	This is a standard complete-information all-pay equilibrium~\citep{baye1996all}, applied to valuations \(G_{i,t}(s)\) and bid cost \(\rho b_i\).
%	The bid-cap assumption ensures that the equilibrium support lies inside \([0,\beta]\).
%	The stated winning probability follows by integrating the probability that \(j^*\)'s bid is below the bid of \(i^*\):
%	\begin{equation}
%	\begin{aligned}
%		P[i^*\text{ wins}]
%		&=
%		\int_0^{G_{j^*,t}(s)/\rho}
%		\left(
%			\frac{G_{i^*,t}(s)-G_{j^*,t}(s)}
%			{G_{i^*,t}(s)}
%			+
%			\frac{\rho b}{G_{i^*,t}(s)}
%		\right)
%		\frac{\rho}{G_{j^*,t}(s)}
%		db\\
%		&=
%		1-\frac{G_{j^*,t}(s)}
%		{2G_{i^*,t}(s)}.
%	\end{aligned}
%	\end{equation}
%\end{proof}



% \section{A Multi-Agent Bidding Approach for Multi-Objective RL}
% \begin{definition}[\momdp]
%     A multi-objective Markov decision process (\momdp) with \(m \in \Z_{>0}\) objectives is specified by a tuple \(\mathcal{M} = (S, A, T, R, \mu)\), where
%     \begin{itemize}
%         \item \(S\) is the set of states,
%         \item \(A\) is the set of actions,
%         \item \(T : S \times A \to \dist(S)\) is the transition function mapping a state-action pair to a distribution over states,
%         \item \(R : S \times A \times S \to \R^m\) is the reward function with each output component corresponding to the different objectives, and
%         \item \(\mu \in \dist(S)\) is the initial state distribution.
%     \end{itemize}
% \end{definition}

% A \textit{policy} in an \momdp~\(\mathcal{M}\) is a function \(\pi : (S \times A)^* \times S \to \dist(A)\) that maps a history of state-action pairs and the current state to a distribution over actions. 

% \begin{definition}[\mabmdp]
%     Let \(\mathcal{M} = (S, A, T, R, \mu)\) be an \momdp~with \(m\) objectives and let \(b \in \Z_{>0}\) be the bid upper bound. Also, define \(M = \{1, \dots, m\}\) be indices of the \(m\) agents corresponding to the \(m\) objectives along with \(\bot\) representing a null agent. Lastly, let \(B = \{0, \dots, b\}\) be the range of bids and \(\rho > 0\) be the bid penalty factor.
%     We define the multi-agent bidding Markov decision process (\mabmdp) as a tuple \(\mathcal{B}_{\mathcal{M}} = (\hat S, \hat{A}, \hat{T}, P, \hat{R}, \hat \mu)\) where
%     \begin{itemize}
%         \item \(\hat S = M \times S\) is the new state space augmented with the index of the agent that won the previous round of bidding,
%         \item \(\hat{A} = A^m \times B^m\) represents the action space of the \(m\) agents in which each agent selects an action from \(A\) and a bid from \(B\),
%         \item \(\hat{T} : \hat S \times \hat{A} \to \dist(\hat S)\) is the new transition function defined as,
%         \begin{equation*}
%             \hat{T}((\_, s), (\vec a, \vec b)) \coloneqq \frac{1}{\abs{B_{\max}}}\sum_{i \in B_{\max}} (T(s, a_i), i)
%         \end{equation*}
%         where \(B_{\max} \coloneqq \{i \mid b_i = \max \{b_1, \dots, b_m\}\}\) is the set of agents with maximal bids. The tuple \((T(s, a_i), i)\) represents the distribution over \(\hat S\) induced by the original transition function \(T\) such that the second component is fixed, and the weighted sum represents taking the weighted sums of the distributions over \(\hat S\).
%         \item \(P : \hat A \times M \to \R^m\) is the bidding penalty for the \(m\) agents and the second component is the index of the agent that won the bidding.
%         \item \(\hat{R} : \hat S \times \hat{A} \times \hat S \to \R^m\) is the reward function for the \(m\) agents with
%         \begin{equation*}
%             \hat{R}_k((\_, s_0), (\vec a, \vec b), (i, s)) \coloneqq R_k(s_0, a_i, s) - P_k((\vec a, \vec b), i)
%         \end{equation*}
%         where \(i \in M\) is the index of the agent that won the bid and chose the action.
%         \item \(\hat \mu \coloneqq (\mu, 1)\) is the initial state distribution over \(\hat S\) induced by \(\mu\) and the second component is fixed to be \(1\) without loss of generality.
%     \end{itemize}
% \end{definition}

% Given an \mabmdp~\(\mathcal{B}_{\mathcal{M}}\), a \textit{policy} for each agent indexed by \(i \in \{1, \dots, m\}\) takes a similar form: \(\pi_i : (\hat S \times \hat A)^* \times \hat S \to \hat A\). Intuitively, a state \((i, s) \in \hat S\) encodes the agent that won the bidding and chose the action to reach \(s\) in the previous step. At each step, each of the agents choose an action and a bid, and an action amongst the set of highest bidders is chosen uniformly at random. The reward function includes a penalty term that captures the desired bidding mechanism.
